\section{Overview}
\label{sec:proposed:overview}

This chapter presents the proposed solution. It is developed in
four steps. Section~\ref{sec:proposed:dataarch} formalises the
data architecture --- both the schema across the GPS and AFC
source systems and the pipeline that produces the
pre-aggregated views the dashboard consumes.
Section~\ref{sec:proposed:eda} reports the exploratory data
analysis on the available HCMC sample, the findings of which
scope every downstream design decision.
Section~\ref{sec:proposed:modules} screens the analytics
modules drafted in the brainstorming phase against the schema
and the EDA findings, separating the modules that the local
data supports from those that require additional data
releases. Section~\ref{sec:proposed:dashboard} specifies the
dashboard: stakeholders and their analytical questions, filter
flexibility, the three-column layout, the canvas-swap
interaction, and the chart catalogue per analytical lens.

The solution rests on the theoretical foundations introduced
in Chapter~\ref{chapter:theory}: the medallion data
architecture, the analytical stack of DuckDB, Polars,
Streamlit, pydeck and Plotly, the three QoS indicators of
Srisakaokul~\textit{et al.}~\cite{srisakaokul2023}, the
FuzzyCut, GroupSet and Gridify visualisation primitives of
Liu~\cite{liu2022}, the visual storytelling principles of
Knaflic~\cite{knaflic2015}, and the Multiple Coordinated Views
interaction model. Where the local sample makes a literature
method only partially applicable (notably the GPS-side QoS
indicators which the local one-day sample limits), the
limitation is documented at the corresponding module rather
than silently bypassed.
